\chapter{Curves and divisors}
\label{mk-ch-curves}

Let \(k\) be a field and let \(C\to\Spec k\) be a smooth, proper,
geometrically integral curve.  Write \(K(C)\) for its function field and
\(\mathcal O_C\) for its structure sheaf.

\begin{definition}[Divisors]
  \label{mk-def-divisor}
  \dcref{II.6,custom}
  \cite{hartshorne-ag,papaioannou-rr}
  A Weil divisor on \(C\) is a finite sum \(D=\sum_{x\in C^{(1)}}n_x[x]\),
  \(n_x\in\mathbb Z\).  Its degree is
  \[
    \deg D=\sum_x n_x[\kappa(x):k].
  \]
  For \(f\in K(C)^\times\), \(\operatorname{div}(f)=\sum_x
  \operatorname{ord}_x(f)[x]\) is principal; divisors modulo principal
  divisors form \(\operatorname{Pic}(C)\).
\end{definition}

\begin{definition}[Genus]
  \label{mk-def-genus}
  \dcref{IV.1,custom}
  \cite{milne-av,kleiman-picard}
  The genus of \(C\) is
  \[
    g(C)=\dim_k H^1(C,\mathcal O_C).
  \]
\end{definition}

\begin{theorem}[Riemann--Roch]
  \label{mk-thm-riemann-roch}
  \dcref{IV.1,custom}
  \uses{mk-def-divisor,mk-def-genus}
  \cite{hartshorne-ag,papaioannou-rr}
  There is a canonical divisor \(K_C\) of degree \(2g(C)-2\) and, for every
  divisor \(D\),
  \[
    \ell(D)-\ell(K_C-D)=\deg D+1-g(C),
    \qquad \ell(D)=\dim_k H^0(C,\mathcal O_C(D)).
  \]
\end{theorem}

\begin{corollary}[Large-degree vanishing]
  \label{mk-cor-large-degree}
  \dcref{IV.1,custom}
  \uses{mk-thm-riemann-roch}
  If \(\deg D>2g(C)-2\), then \(H^1(C,\mathcal O_C(D))=0\) and
  \(\ell(D)=\deg D+1-g(C)\).
\end{corollary}
